\chapter{System Architecture}
\label{ch:system-architecture}
% Port from: docs/04_system_architecture.md

\section{Overview}
\projectname{} is a two-tier web system. The client is a React~19
single-page application (Vite, TypeScript, TailwindCSS~v4) that talks to the
server over two channels: REST for all resource operations, via RTK~Query
API slices with declarative cache invalidation, and a WebSocket for the
real-time navigation stream. The server is built on
\code{express-file-cluster} (EFC), a file-system-routed Express framework:
every file under \code{src/api/} maps to a URL path, so adding an endpoint
means adding a file. The server runs as a Node.js cluster (one primary, two
workers sharing a port), persists to MongoDB via Mongoose, queues
transactional email through a Redis-backed BullMQ queue, and serves waypoint
media from Cloudinary.

\begin{figure}[H]\centering
  \includegraphics[width=\textwidth]{architecture}
  \caption{System architecture: client, server, and data layers.}\label{fig:architecture}
\end{figure}


\section{Client}
Routing is role-aware: a \code{ProtectedRoute} component reads the
authenticated identity and confines regular users to the exploration
experience (dashboard, explore, people, redeem, profile) while admins are
routed to a separate dashboard, user management, and the waypoint logger ---
neither role can reach the other's surface. Data fetching is centralised in
RTK~Query slices (auth, places, routes, progress, admin, points/quiz), so a
mutation such as confirming a visit or redeeming a reward declaratively
invalidates exactly the caches that depend on it. The AR layer (A-Frame and
AR.js) is deliberately kept outside the React tree and is loaded only when
navigation starts; Chapter~\ref{ch:ar-pathfinder} covers it in detail.

\section{Server}
REST requests pass through cookie-based JWT authentication
(\code{requireAuth} middleware attaches \code{\{id, role, email\}} to the
request), then reach file-routed handlers that call the models. Business
logic that spans endpoints lives in a small service layer:
\code{proximity.ts} (distance, bearing, arrival, confirmation),
\code{arPoints.ts} (award/spend ledger), and \code{tripSummaries.ts}
(per-place progress and badge computation).

Authentication uses short-lived access JWTs and longer-lived refresh JWTs,
both in HTTP-only cookies so tokens are invisible to page JavaScript;
passwords are stored as bcrypt hashes. The WebSocket cannot carry cookies
portably, so it uses a ticket pattern: the client requests a short-lived,
single-purpose JWT over authenticated REST and presents it as a query
parameter; the server verifies it during the HTTP upgrade, before the socket
is accepted.

\section{Real-Time Location Pipeline}
Navigation is a server-authoritative loop:

\begin{enumerate}
  \item The client seeds an idempotent progress record
        (\apiroute{POST}{/api/user-progress}), fetches a WS ticket, and
        connects to \code{/ws/pathfinder}.
  \item \code{navigator.geolocation.watchPosition} streams coordinates;
        each fix is forwarded as a \code{location} frame.
  \item For every frame the server finds the next unvisited waypoint,
        computes haversine distance and bearing, and replies with a
        \code{progress} frame carrying the distance, bearing, the place's
        geofence radius (\code{visit\_threshold\_meters}, default 10\,m,
        stored per place and read live from the database), and an
        \code{arrived} flag.
  \item Arrival never auto-completes a checkpoint. The client shows a
        confirmation prompt; only an explicit \code{confirm\_visit} message
        --- re-validated server-side against the last reported position and
        the same radius --- marks the waypoint visited.
  \item Confirmation triggers the gamification side effects
        (Chapter~\ref{ch:gamification}): milestone and side-quest point
        awards, the permanent milestone log on the user profile, and, when
        the final checkpoint falls, the place-completion award plus a
        one-time increment of the place's visitor counter.
\end{enumerate}

Because the threshold, the distance check, and the awards all live
server-side, a confirmed visit constitutes evidence of physical presence
that a client cannot forge by manipulating its own state.

\begin{figure}[H]\centering
  \includegraphics[width=\textwidth]{location-pipeline}
  \caption{Real-time location pipeline: from GPS fix to checkpoint confirmation.}\label{fig:pipeline}
\end{figure}
